\section*{Introduction: structural determination and information recovery}
\addcontentsline{toc}{section}{Introduction: structural determination and information recovery}
\label{libro:prologo}
The holographic extreme and mean ratio addresses a precise question: how a structure can produce a numerical evaluation, undergo a transformation, and preserve the information required to recognize and invert that transformation. The problem concerns both the formation of the value and the persistence of its relations. An equality between quantities establishes a balance; a recovery law additionally determines which state realizes it, which data identify that state, and which operation restores it. This work develops the latter requirement from an arithmetic construction of paths through to its linear, incidence and physical representations.

The first result organizing the argument is the determination of the holographic arithmetic--geometric coupling constant. Its ordered representation is the dodecaphasic register \(K\); its positional evaluation is the rational number \(\kappa\). This reading preserves the register exactly when the base, length and phase origin are fixed. The object is also significant through its operations: its cyclic shifts form a basis, its responses permit operator identification, and its incidence images transport the relations of the initial configuration. Structure, equation and value thus arise as stages of a construction whose recoverable information is determined mathematically.

The second result supplies the general law of that conservation. Two isometric transports with a common domain and codomain produce two components whose weighted balance restores the initial inner product exactly. The nonadic specialization retains a ternary norm of value \(73\) and determines an observable transformation with weights \(72/73\) and \(1/73\). A policy of prolongation and archiving turns this local identity into stable recovery at arbitrary depth. The article's thesis thereby acquires a verifiable content: the evaluated value, the preserved state and the observed readout are related by explicit maps and inverse formulae.

\subsection*{The generating kernel and its two realizations}
The construction is established in a formal kernel called \emph{Triadic Modular Holography}, hereafter HMT. Its first object is \emph{All Possible Paths}, hereafter APP: an arithmetic path structure on the toroidal support \((\mathbb Z/9\mathbb Z)^2\), whose Cayley graph describes cardinal displacements. Additive and multiplicative evaluations are distinguished on this same support. Preserving the residue together with the quotient retains the evaluated integer and permits the origin of each evaluation to be traced. Addition accumulates; multiplication organizes the second sheet through its own evaluation and transport rules.

The second object, \emph{TRIT}, supplies the oriented signature \(\{-1,0,+1\}\) and distinguishes the local regimes. The extended Euler equation expresses the elliptic, parabolic and hyperbolic regimes in their quadratic algebras through a common exponential law. The third object is the \emph{Topological Phase Kernel}, hereafter TPK. Its action composes selection, transport and memory updating. Route, direction, sheet, residue, quotient, carry and boundary are retained in the relevant enriched state; phase return prolongs its history. The nonadic connection and holonomy belong to this dynamics.

This state gives rise to the joint discrete structure of the continuum and its two realizations: contraction of compatible cylinders determines arithmetic coordinates, while boundary incidence preserves combinatorial and geometric relations. The five correlative constructions of the continuum remain within this same development. The linear spaces, forms and observables used subsequently realize previously generated objects, with their domains and normalizations. The continuity of the argument lies in each operation preserving the relations required by the next.

\subsection*{Completion of unity and reversible refinement}
Completing a cell with nine possibilities per coordinate adds one distinguished boundary position. In three coordinates, stratification by the number of boundary positions gives
\[
10^3=729+243+27+1=729+271.
\]
The interior, intermediate strata and apex compose one normalized unity. In arbitrary finite dimension, the expansion of \((9+1)^d\) expresses the same stratification; summing the contributions of the omitted coordinate proves compatibility with projection.

The central block \(\bigl(\begin{smallmatrix}7&2\\2&7\end{smallmatrix}\bigr)\) provides a concrete arithmetic antecedent. Its eigenvalues are \(9\) and \(5\), and its concatenations satisfy \(72+27=77+22=99\). Restoring one unit and then refining gives the chain \((72,27)\mapsto(73,27)\mapsto(729,271)\). The first operation incorporates the recorded unit; the second changes scale and redistributes one unit by a zero-sum correction. This explicit sequence relates centre, completion and proportion while retaining the origin of every term.

Refinement with carry introduces an evolution of this distribution. The transformation \((x,y)\mapsto(Bx-c,By+c)\) preserves the sum after multiplication by \(B\). On admissible integer labels, with \(0\leq c<B\), Euclidean division of the second component first recovers \(c\) and then the preceding pair. Depth fixes the number of inverse steps. When lifted carries are admitted, their additional quotient belongs to the preserved register. Reversibility therefore depends on precise data and has an explicit recovery operation.

This distinction gives holographic conservation an operational meaning. The normalized reading determines a coordinate; the register preserves the relations permitting state reconstruction or further refinement. The subsequent development makes this correspondence linear: the bijection between admissible labels induces a unitary operator between their bases and permits amplitudes, correlations and observables to be transported.

\subsection*{The arithmetic and incidence roles of \(K\)}
The register \(K\) is obtained from the terminal state through its visits, traces, counts, differences and charge. The evaluation
\[
\kappa=\frac{N_K}{1000^{12}-1}
\]
is reversible in the fixed chart: multiplication by the denominator restores the integer, and its decomposition into twelve blocks recovers the complete register, including leading zeroes. The enriched history additionally carries the depth and transport coordinates used by each prolongation. Recovery of each object is formulated on the data that identify it.

In the arithmetic reading, \(\kappa\) enters
\[
\alpha_A=\pi_{\rm HMT}+e_{\rm HMT}-\varphi_{\rm HMT}-4-\kappa.
\]
The three regional coordinates and the register arise from the same generating kernel and act in their constructive order. The complete analytic chart of \(\alpha\) represents this same pre-coordinate through a simple root and a compatible interval collection. In the incidence reading, the register transports its phase relations to the exceptional configuration. The unity of the object is manifested in the effective composition of its readers, inverses and actions.

The cyclic orbit of \(K\) adds an identification result. Its twelve vectors form a basis of the twelve-component linear realization. The responses to this basis determine any linear operator in the class considered; for an operator commuting with the shift, one complete vector response determines the others. The smallest Gram eigenvalue, \(589609\), provides an exact stability bound. Two complementary memory observations also recover the centred part of the register; the uniform charge completes its reconstruction. The constant thus functions as a coupling object with quantified capacities for evaluation, transport and identification.

\subsection*{Bilateral balance and observable transformation}
The integer \(73\) occurs as the norm of two conjugate elements and as the index of a rank-two lattice inclusion. The identity \(10\cdot73=729+1\) relates this ternary norm to decimal completion. The bilateral theorem extends the balance to two isometries \(B,C\) between spaces with a common domain and codomain. For \(a>0\), the components \(aB-C\) and \((a+1)B+C\), weighted by \(a+1\) and \(a\), preserve the inner product after the stated normalization. Expansion and cancellation of the cross terms prove the identity in arbitrary dimension. The specialization \(a=8\) retains the nonadic operators and their arithmetic information.

Conservation of the complete state permits the effect of a particular observation to be calculated. If the relative transport \(R\) is unitary and \(G\) is the observable expressed in the common chart, the diagonal reading transforms
\[
\mathcal T(G)=\frac{72}{73}G+\frac1{73}R^*GR.
\]
An observable remains invariant exactly when it commutes with \(R\). The observable transported on the complete space preserves the entire initial reading. Within a single construction, the result specifies both the information that persists and the change introduced by observing it through a particular channel.

\subsection*{Recovery at arbitrary depth}
At each step, one component may be prolonged and its complement archived. The finite balance sums the squared norms of the memories and the terminal state. An explicit policy imposes the geometric bound \((729/1241)^N\) on the squared norm of the nonadic terminal operator. In the limit, the entire norm is contained in the infinite archive; this archive is an isometry and its adjoint reconstructs the state with operator norm one.

The first archived blocks likewise admit an exact inverse corrected by their Gram operator, with a stability bound uniform over positive depths. The infinite archive, the finite inverse and the approximation obtained by truncating the adjoint have their own formulae and estimates. In general dynamics with a surviving terminal, its contribution is retained together with memory. The proved policy identifies the case in which that terminal tends to zero.

The unitary lift of refinement composes with bilateral balance. The decimal example carrying \((73,27,1)\) to \((729,271)\) keeps the origin of the carry explicit; comparison with a neighbouring carry determines the exact change in a reading. The evolutionary conservation of unity is thereby expressed through a sequence of recoverable transformations.

\begin{figure}[htbp]
\centering
\begin{tikzpicture}[>=Stealth,every node/.style={font=\small,align=center},
block/.style={draw=ink,rounded corners=2pt,inner sep=6pt,text width=35mm},
wide/.style={draw=ink,inner sep=6pt,text width=47mm}]
\node[block] (d) {Admissible arithmetic\\state\\\((x,y,c)\)};
\node[block,right=14mm of d] (e) {Refined state\\\((10x-c,10y+c)\)};
\node[block,right=14mm of e] (j) {Label bases\\\(J\) unitary};
\draw[<->] (d)--node[above,font=\scriptsize]{carry} (e);
\draw[->] (e)--(j);
\node[wide,below=17mm of e] (w) {Bilateral transport\\\(W^*W=I\)};
\draw[->] (j.south)|-(w.east);
\node[block,below left=17mm and 4mm of w] (a) {Joint reading\\\(\frac{72G+R^*GR}{73}\)};
\node[block,below right=17mm and 4mm of w] (b) {Transported observable\\\(WGW^*\)};
\draw[->] (w)--(a);\draw[->] (w)--(b);
\node[font=\small,below=8mm of a] {Invariant if \(GR=RG\)};
\node[font=\small,below=8mm of b] {Recovery through \(W^*\)};
\end{tikzpicture}
\caption{Composition of proved operations. Refinement preserves the normalized sum and allows the label to be recovered. Its lift acts on state amplitudes and preserves the norm. Reading an observable and transporting it in full retain different functions, specified by the two lower formulas.}
\label{libro:fig:composicion}
\end{figure}


\subsection*{Forms, action and quantum observation}
Inner-product transport extends to every admissible exterior degree. Areas and volumes also preserve the mixed readout--memory contributions. Exceptional incidence transports these identities on its specified images; variational charges are obtained from a declared action and its symmetries. The electromagnetic and gravitational realizations retain the scales and operators involved in their proofs.

The Gaussian realization shows why preserving the state and knowing a reduced observation are complementary questions. In the declared protocol with two identical centred pure Gaussian inputs and finitely many modes, symplectic transport preserves the global structure, while its reduction acquires a spectrum determined by the antilinear component of memory. The elliptic--parabolic loop studied here produces purity \(9/11\), spectrum \((9/10)10^{-j}\), and dimensionless entropy \((10/9)\log10-\log9\), with all Fock-space occupations retained in the calculation.

The argument therefore concludes with a transformation law richer than a scalar equality: the value is determined, the structure involved in its reading is preserved, and the conditions for its recovery are calculated. The inverse formulae, stability bounds and specified realizations jointly express the mathematical meaning of the holographic extreme and mean ratio.

% BEGIN SINTESIS_ADITIVA_20260922
\par
Section~\ref{delta22:xi:holonomia} makes the archive contribution to bilateral composition explicit, constructs the quarter-turn plane of the dodecaphasic calendar, and proves oriented action recovery through a relative commutator. Compatible amplification preserves the operator; the detector distinguishes state dilation from energy gain and the ordered archive from operation counting.
% END SINTESIS_ADITIVA_20260922
